% GENERATED FILE. Edit canonical lessons, then run npm run generate:books.
% canonical-source-hash: fnv1a64:38898a2ac6aba29a
% canonical-lessons: RU-C469-revnost, RU-C469-revnovat, RU-C469-revnivyy, RU-C469-zavist, RU-C469-zavidovat

\chapter{Jealousy and Envy}
\label{ch:ru-jealousy-and-envy}
\hlchaptermodality{469}

\begin{chapteropening}
\textbf{By the end of this chapter:} \emph{I can talk about jealousy and envy.}

Complete the last lesson of chapter 469: I can talk about jealousy and envy.
\end{chapteropening}

\section[\texorpdfstring{\ru{ревность} (révnost)}{ревность (révnost)}]{\ru{ревность} (révnost) — jealousy}

\label{lesson:RU-C469-revnost}



\begin{quote}
Before the new one: say the Russian for a task, then the Russian for experience.
\end{quote}

\subsection*{You'll want to know: \ru{ревность}}
\textbf{\ru{ревность}} (\emph{révnost}) — \textquotedblleft{}jealousy\textquotedblright{}.

\begin{quote}
\ru{Из}-\ru{за} \ru{ревности} \ru{они} \ru{часто} \ru{ссорятся}.
\end{quote}

\begin{quote}
\emph{Iz-za révnosti oní chásto ssóryatsya.}
\end{quote}

\begin{quote}
They often quarrel because of jealousy.
\end{quote}

Feminine; no plural in everyday use.

\begin{cousinweb}
\emph{\ru{ревность}} is a noun in \emph{-\ru{ость}} on the stem \emph{\ru{ревн}-}, which it shares with \emph{\ru{ревновать}} and \emph{\ru{ревнивый}}; EtymDB links it to Church Slavonic \emph{rĭvenie} 'strife, quarrel'. An older sense, 'zeal', is still met in bookish Russian.
\end{cousinweb}

\begin{grammarlens}[title={What you've built}]
One more word for telling, explaining and describing at B1.
\end{grammarlens}

\subsection*{Your turn}
\begin{itemize}
\raggedright
  \item \emph{Say it:} \emph{révnost}
  \item \emph{Say it:} \emph{révnost}, once more
  \item \emph{Say it:} a sentence of your own with \emph{révnost}
  \item \emph{Recall:} say the Russian for a task, then the Russian for experience, then say \emph{révnost} again
\end{itemize}

\begin{tcolorbox}[breakable,colback=teal!4,colframe=teal!35!black,title={Before you move on}]
What does \ru{ревность} mean? (\textquotedblleft{}Jealousy\textquotedblright{}.) Say it once more.
\end{tcolorbox}
\section[\texorpdfstring{\ru{ревновать} (revnovát)}{ревновать (revnovát)}]{\ru{ревновать} (revnovát) — to be jealous}

\label{lesson:RU-C469-revnovat}



\begin{quote}
Before the new one: say the Russian for experience, then the Russian for jealousy.
\end{quote}

\subsection*{You'll want to know: \ru{ревновать}}
\textbf{\ru{ревновать}} (\emph{revnovát}) — \textquotedblleft{}to be jealous\textquotedblright{}.

\begin{quote}
\ru{Она} \ru{ревнует} \ru{мужа} \ru{к} \ru{его} \ru{коллегам}.
\end{quote}

\begin{quote}
\emph{Oná revnúyet múzha k yegó kollégam.}
\end{quote}

\begin{quote}
She gets jealous of her husband's colleagues.
\end{quote}

Imperfective; the present is \ru{ревную}, \ru{ревнуешь} (-\ru{ова}- becomes -\ru{у}-). \emph{\ru{Ревновать} \ru{кого} \ru{к} \ru{кому}} — to be jealous about someone because of someone. Perfective \ru{приревновать}.

\begin{cousinweb}
\emph{\ru{ревновать}} is built with \emph{-\ru{овать}} on the stem it shares with \emph{\ru{ревность}} 'jealousy'; EtymDB lists the verb without a source and links the noun to Church Slavonic \emph{rĭvenie} 'strife'.
\end{cousinweb}

\begin{grammarlens}[title={What you've built}]
One more word for telling, explaining and describing at B1.
\end{grammarlens}

\subsection*{Your turn}
\begin{itemize}
\raggedright
  \item \emph{Say it:} \emph{revnovát}
  \item \emph{Say it:} \emph{revnovát}, once more
  \item \emph{Say it:} a sentence of your own with \emph{revnovát}
  \item \emph{Recall:} say the Russian for experience, then the Russian for jealousy, then say \emph{revnovát} again
\end{itemize}

\begin{tcolorbox}[breakable,colback=teal!4,colframe=teal!35!black,title={Before you move on}]
What does \ru{ревновать} mean? (\textquotedblleft{}To be jealous\textquotedblright{}.) Say it once more.
\end{tcolorbox}
\section[\texorpdfstring{\ru{ревнивый} (revnívyy)}{ревнивый (revnívyy)}]{\ru{ревнивый} (revnívyy) — jealous (by nature)}

\label{lesson:RU-C469-revnivyy}



\begin{quote}
Before the new one: say the Russian for jealousy, then the Russian for to be jealous.
\end{quote}

\subsection*{You'll want to know: \ru{ревнивый}}
\textbf{\ru{ревнивый}} (\emph{revnívyy}) — \textquotedblleft{}jealous (by nature)\textquotedblright{}.

\begin{quote}
\ru{Мой} \ru{муж} \ru{очень} \ru{ревнивый}.
\end{quote}

\begin{quote}
\emph{Moy muzh óchen revnívyy.}
\end{quote}

\begin{quote}
My husband is very jealous.
\end{quote}

Describes a jealous character. Short forms: \ru{ревнив}, \ru{ревнива}.

\begin{cousinweb}
\emph{\ru{ревнивый}} is built on the stem of \emph{\ru{ревность}} 'jealousy' with the adjective ending \emph{-\ru{ивый}}; EtymDB lists the adjective without a source.
\end{cousinweb}

\begin{grammarlens}[title={What you've built}]
One more word for telling, explaining and describing at B1.
\end{grammarlens}

\subsection*{Your turn}
\begin{itemize}
\raggedright
  \item \emph{Say it:} \emph{revnívyy}
  \item \emph{Say it:} \emph{revnívyy}, once more
  \item \emph{Say it:} a sentence of your own with \emph{revnívyy}
  \item \emph{Recall:} say the Russian for jealousy, then the Russian for to be jealous, then say \emph{revnívyy} again
\end{itemize}

\begin{tcolorbox}[breakable,colback=teal!4,colframe=teal!35!black,title={Before you move on}]
What does \ru{ревнивый} mean? (\textquotedblleft{}Jealous (by nature)\textquotedblright{}.) Say it once more.
\end{tcolorbox}
\section[\texorpdfstring{\ru{зависть} (závist)}{зависть (závist)}]{\ru{зависть} (závist) — envy}

\label{lesson:RU-C469-zavist}



\begin{quote}
Before the new one: say the Russian for to be jealous, then the Russian for jealous.
\end{quote}

\subsection*{You'll want to know: \ru{зависть}}
\textbf{\ru{зависть}} (\emph{závist}) — \textquotedblleft{}envy\textquotedblright{}.

\begin{quote}
\ru{Коллеги} \ru{говорили} \ru{о} \ru{его} \ru{успехе} \ru{без} \ru{зависти}.
\end{quote}

\begin{quote}
\emph{Kollégi govoríli o yegó uspékhe bez závisti.}
\end{quote}

\begin{quote}
His colleagues talked about his success without envy.
\end{quote}

Feminine. \emph{\ru{С} \ru{завистью}} — enviously.

\begin{cousinweb}
\emph{\ru{зависть}} is built on \emph{\ru{за}-} and the root of \emph{\ru{видеть}} 'to see', so envy is a way of looking at what another person has; \emph{\ru{видеть}} is inherited from Proto-Slavic \emph{viděti}, from PIE \emph{weyd- 'to see, to know' (EtymDB).}
\end{cousinweb}

\begin{grammarlens}[title={What you've built}]
One more word for telling, explaining and describing at B1.
\end{grammarlens}

\subsection*{Your turn}
\begin{itemize}
\raggedright
  \item \emph{Say it:} \emph{závist}
  \item \emph{Say it:} \emph{závist}, once more
  \item \emph{Say it:} a sentence of your own with \emph{závist}
  \item \emph{Recall:} say the Russian for to be jealous, then the Russian for jealous, then say \emph{závist} again
\end{itemize}

\begin{tcolorbox}[breakable,colback=teal!4,colframe=teal!35!black,title={Before you move on}]
What does \ru{зависть} mean? (\textquotedblleft{}Envy\textquotedblright{}.) Say it once more.
\end{tcolorbox}
\section[\texorpdfstring{\ru{завидовать} (zavídovat)}{завидовать (zavídovat)}]{\ru{завидовать} (zavídovat) — to envy}

\label{lesson:RU-C469-zavidovat}



\begin{quote}
Before the new one: say the Russian for jealous, then the Russian for envy.
\end{quote}

\subsection*{You'll want to know: \ru{завидовать}}
\textbf{\ru{завидовать}} (\emph{zavídovat}) — \textquotedblleft{}to envy\textquotedblright{}.

\begin{quote}
\ru{Я} \ru{не} \ru{завидую} \ru{богатым} \ru{людям}.
\end{quote}

\begin{quote}
\emph{Ya ne zavíduyu bogátym lyúdyam.}
\end{quote}

\begin{quote}
I don't envy rich people.
\end{quote}

Imperfective; the present is \ru{завидую}, \ru{завидуешь}. Takes the dative: \emph{\ru{завидовать} \ru{брату}}. Perfective \ru{позавидовать}.

\begin{cousinweb}
\emph{\ru{завидовать}} is built on the same \emph{\ru{за}-} and \emph{\ru{вид}-} as \emph{\ru{зависть}} 'envy'; the root is that of \emph{\ru{видеть}} 'to see', inherited from Proto-Slavic \emph{viděti} (EtymDB).
\end{cousinweb}

\begin{grammarlens}[title={What you've built}]
One more word for telling, explaining and describing at B1.
\end{grammarlens}

\subsection*{Your turn}
\begin{itemize}
\raggedright
  \item \emph{Say it:} \emph{zavídovat}
  \item \emph{Say it:} \emph{zavídovat}, once more
  \item \emph{Say it:} a sentence of your own with \emph{zavídovat}
  \item \emph{Recall:} say the Russian for jealous, then the Russian for envy, then say \emph{zavídovat} again
\end{itemize}

\begin{tcolorbox}[breakable,colback=teal!4,colframe=teal!35!black,title={Before you move on}]
What does \ru{завидовать} mean? (\textquotedblleft{}To envy\textquotedblright{}.) Say it once more.
\end{tcolorbox}
